% Chapter 9 — Transactional Workflows and API Mapping

\chapter{Transactional Workflows \& API Mapping}
\label{ch:tx}

\section{What we mean by ``transaction'' here}
In everyday speech, a transaction is ``a complete action.'' In the database sense, it can be a \texttt{BEGIN ... COMMIT} block; in this project many operations use Supabase single-row inserts that still behave atomically per statement. The important idea for reports is \textbf{order of checks}: validate permissions first, then write facts that constraints protect.

\section{Enrollment workflow (student joins a course)}
\begin{enumerate}[leftmargin=*]
  \item Load the course row. Confirm it is published (for students) and not soft-deleted.

  \item Check whether an enrollment already exists for this (\texttt{user\_id}, \texttt{course\_id}). If yes, return a friendly error instead of hitting the UNIQUE constraint.

  \item Resolve the ``active'' status ID from \texttt{enrollment\_statuses}.

  \item Insert into \texttt{enrollments}.
\end{enumerate}

\textbf{Why uniqueness matters:} Two tabs open at once might both try to enroll. The application may race; the database UNIQUE constraint still prevents duplicate rows---correctness does not depend only on JavaScript logic.

\section{Quiz submission workflow}
\begin{enumerate}[leftmargin=*]
  \item Confirm the student may take this quiz: enrolled (where required), quiz published, optional single-attempt rule enforced by checking existing \texttt{quiz\_attempts}.

  \item Load all questions and options for scoring.

  \item Grade server-side: compare chosen \texttt{option\_id} with \texttt{is\_correct}; weight by question points; compute percentage score; compare with \texttt{pass\_score}.

  \item Insert one \texttt{quiz\_attempts} row, then insert many \texttt{quiz\_answers} rows.

\end{enumerate}

\textbf{Why grade on server:} The browser cannot be trusted to calculate scores---users could edit DevTools. The database stores the authoritative attempt.

\section{Mapping endpoints to tables (examples)}
\begin{itemize}[leftmargin=*]
  \item \texttt{POST /courses/:id/content} $\rightarrow$ insert \texttt{content} (and resolve \texttt{content\_types} by name).

  \item \texttt{POST /activity} $\rightarrow$ resolve \texttt{activity\_types} by name, insert \texttt{activity\_log}.

  \item \texttt{PUT /progress/:contentId} $\rightarrow$ upsert \texttt{content\_progress}.

  \item \texttt{GET /analytics/...} $\rightarrow$ read \texttt{mv\_*} views and/or aggregate \texttt{activity\_log} / \texttt{quiz\_attempts} depending on role and filters (instructor paths often recompute with enrollment filters).
\end{itemize}

\section{Concurrency in one sentence}
``We rely on constraints like UNIQUE for enrollments so two simultaneous requests cannot corrupt state even if the server code is imperfect.''
